微网通过在外部电网电价低谷或分布式光伏盈余时段为储能充电、在电价高峰时段放电，可在满足小区负载的前提下降低购电费用。本文建立\textbf{统一的逐时段确定性线性规划框架}，把时间分辨率、信息结构、结算制度与电价情景四个维度分别参数化，依次求解四个问题，交付 \texttt{result1.xlsx}、\texttt{result2.xlsx}、\texttt{result3.xlsx}、\texttt{result4-2.xlsx}、\texttt{result4-3.xlsx} 五个结果文件。

\textbf{模型与思路。} 以 10 分钟为决策时段，决策变量为计划购电量、调整购电量、紧急购电量、充放电量与弃光量，状态量为储电量。功率平衡取等式，储能动态计入 90\% 双向效率（往返 81\%），容量与功率约束按题面附录。目标为购电费用最小且\textbf{不乘时段长}。本文给出互补性引理：往返效率小于 1 且电价为正时，“同时充放电”被“少充少放”严格支配，故线性松弛无损、无需二元变量；鉴于最优调度面可能不唯一，全部模型统一按固定规则选取输出的解。

\textbf{主要结果。} 问题一在附件 1 的代表性单日上得最优全天购电费 $35126.948589$ 元、购电量 $59482.698998$ kWh，呈典型的分时电价套利形态，最便宜的 20\% 电价时段承担 55.4\% 的充电量，日周期精确闭合、全日无弃光。问题二把模型滚动到全年 365 天（52\,560 时段），得全期购电费 $13\,758\,182.573724$ 元、交付期（2 月 1 日—12 月 31 日，334 天）$12\,233\,050.830708$ 元，并证明在“完全信息 + 无购电上限”下紧急购电恒为零。问题三引入附件 3 整点光伏预报的降尺度与 6:00、12:00、18:00 三个调整时刻及 $0.5\times/1.5\times$ 双向分段结算，交付期费用升至 $14\,540\,616.335652$ 元，比问题二高 $2\,307\,565.50$ 元（$+18.863\%$），该差额即预报误差的代价。问题四在附件 4 的实时波动电价下重算问题二与问题三，两条链的交付期费用分别为 $13\,006\,411.041148$ 元与 $15\,289\,050.714738$ 元；两链差额 $2\,282\,639.673590$ 元可\textbf{完全归因}为计划量差异 $100\,230.446225$ 元、偏差结算 $522\,512.788776$ 元与紧急购电 $1\,659\,896.438588$ 元三项。按“0:00 时当日电价未知”的处理方式，决策层使用由历史价格构造的预测价、结算一律用实际价，价格预报误差的代价为 $520\,790.839035$ 元与 $258\,485.465704$ 元。

\textbf{验证与稳健性。} 四问均通过五项独立检验：结果完整性、144 时段独立回代、量纲与逐层目标复算、解析界与常识检验、跨问题一致性五级验收，等式残差不超过 $4.6\times10^{-12}$，交付工作簿逐格复算一致。问题一判“稳定”，问题二判“脆弱（需缩小适用范围）”，问题三与链 \emph{4-3} 判“条件稳定”，链 \emph{4-2} 判“稳定”；全部未通过的检验与反例均如实列出，未通过调参或放宽阈值使其“通过”。

\textbf{创新点。} 以四个维度上的参数化统一四问并形成可复用模板；用互补性引理把“是否需要二元变量”从工程判断变为数学结论；把不确定性代价拆成可计价的分项；用退化情景复现构造实现正确性基准（六项指标逐位复现，相对差不超过 $4.86\times10^{-13}$）；用只改“看多远”的滚动时域对照回答结构性问题；并对“是否需要更多预报、是否需要随机优化”给出三层递进的定量回答。不足在于文献证据仅达摘要/元数据级、未显式强制充放电互补、若干物理简化，以及全部结果基于确定性点预测，已在正文逐条说明。
